\section{Domain-Adapted Language Models}

Domain adaptation has proven to be a major leverage point for enhancing the performance of embedding models in specialised fields \cite{gururangan_dont_2020}. At its core, this adaptation relies on fine-tuning, a form of transfer learning, which \citet{Jurafsky2024_SLP} define as the process of tailoring a pre-trained language model to a specific domain or task. Models are first trained on large, generic corpora to learn foundational linguistic structures and subsequently fine-tuned on domain-specific data to capture the unique characteristics of the target domain \cite{gururangan_dont_2020}. This approach makes it possible to bridge the gap between the distributions of the training corpora and the target domain data, which has been determined to be the cause of the poor performance of standard embedding methods \cite{Lee_2019_BioBERT, beltragy_2019_scibert,gururangan_dont_2020,chalkidis-etal-2020-legal}.

Therefore domain specific models have been developed across several fields. In the biomedical domain, Bio\ac{BERT} by \citet{Lee_2019_BioBERT} and Sci\ac{BERT} by \citet{beltragy_2019_scibert} demonstrated the efficacy of domain-adaptive pre-training through further pre-training on domain corpora. Clinical \ac{BERT}  by \citet{alsentzer_2019_clinical_bert} extended this to clinical notes. In the financial domain, Fin\ac{BERT} by \citet{yang_finbert_2020} adapted \ac{BERT}  for financial communications, while BloombergGPT by \citet{wu2023bloomberggpt} later demonstrated domain adaptation at the generative model scale.

Within the legal domain, \citet{chalkidis-etal-2020-legal} provided the most systematic investigation of domain adaptation strategies for \ac{BERT}, testing three approaches. First using \ac{BERT}  out of the box, second, further fine-tuning on domain-specific corpora and third, training from scratch with a domain-specific corpus-vocabulary. Their resulting model family, LEGAL-\ac{BERT} , was trained on 12 GB of diverse English legal text and consistently outperformed generic baselines on downstream legal tasks, such as text classification on EURLEX57K, judgment prediction on ECHR cases and named entity recognition on contracts.  

Furthermore, \citet{hua2022legalrelectra} explored with LegalReLECTRA architectural adjustments. They combined legal and medical corpora with a REFORMER-based architecture to process passages of up to 8,092 tokens. This addressed a practical limitation of standard \ac{BERT} models whose 512-token window is insufficient for lengthy legal documents. evaluations have remained focused on classification and sequence-labelling tasks. Moreover, \citet{niklaus_2024_multilegalpile} created with MultiLegalPile an 689 GB multilingual legal corpus spanning 24 languages, wich he then used to pre-train a series of legal language models that established new state-of-the-art results on LexGLUE and LEXTREME, two comprehensive legal \ac{NLP} benchmarks \cite{niklaus_2024_multilegalpile,chalkidis_2022_lexglue, niklaus_2023_lextreme}\footnote{Due to the rapid proliferation of domain-adapted language models across multiple fields, an exhaustive enumeration of all available embedding architectures is impractical. Extensive reviews of legal-specific models are provided by \citet{chalkidis_2022_lexglue}.}.

While MultiLegalPile includes German-language legal data from Swiss federal court rulings, it does not address the German tax law domain specifically, nor does it evaluate embedding representations for retrieval tasks.

The German Tax Law Domain is directly addressed by two recent lines of work. \citet{wind2026steuerllm} presented SteuerLLM, a domain-adapted large language model with 28 billion parameters specifically for German tax law, trained on synthetically generated data from authentic examination material via a controlled \ac{RAG} pipeline. Alongside the model, they released SteuerEx, the first open benchmark for German tax law, comprising 115 questions drawn from undergraduate and graduate university tax law examinations across six core domains. SteuerLLM demonstrates that a domain-adapted model outperforms substantially larger general-purpose models on realistic tax law reasoning tasks, confirming that domain-specific adaptation yields measurable benefits in this domain \cite{wind2026steuerllm}. The work of \citet{luketina2025using} takes a complementary approach and compares \ac{RAG} architectures with the fine-tuning of GPT-4o specifically for Austrian VAT law. Crucially, both \citet{wind2026steuerllm} and \citet{luketina2025using} address domain adaptation exclusively at the generative model level, optimising the \ac{LLM} that formulates the final answer, while treating the retrieval component as a fixed, underlying pipeline element. 
